% Version 5, 12 November 2026. Revised from v004.tex.
% Bibliography: references.bib, shared with the companion paper gaga_kahler_lefschetz.tex.

\documentclass[11pt,reqno]{amsart}

% ----------------------------------------------------------------------
% Packages
% ----------------------------------------------------------------------
\usepackage[T1]{fontenc}
\usepackage{lmodern}
\usepackage{amssymb}
\usepackage{microtype}
\usepackage[a4paper,margin=1.15in]{geometry}
\usepackage{booktabs,tabularx}
\usepackage{xcolor}
\usepackage{hyperref}

\hypersetup{
  colorlinks=true,
  linkcolor=blue!60!black,
  citecolor=blue!60!black,
  urlcolor=blue!60!black,
  pdftitle={Notes on Hodge Classes and Algebraic Cycles}
}

% ----------------------------------------------------------------------
% Theorem environments
% ----------------------------------------------------------------------
\newtheorem{theorem}{Theorem}[section]
\newtheorem{proposition}[theorem]{Proposition}
\newtheorem{conjecture}[theorem]{Conjecture}

\theoremstyle{definition}
\newtheorem{definition}[theorem]{Definition}
\newtheorem{example}[theorem]{Example}

\theoremstyle{remark}
\newtheorem{remark}[theorem]{Remark}

\numberwithin{equation}{section}

% ----------------------------------------------------------------------
% Commands
% ----------------------------------------------------------------------
\newcommand{\C}{\mathbb{C}}
\newcommand{\R}{\mathbb{R}}
\newcommand{\Z}{\mathbb{Z}}
\newcommand{\Q}{\mathbb{Q}}
\newcommand{\PP}{\mathbb{P}}
\newcommand{\OO}{\mathcal{O}}
\newcommand{\an}{\mathrm{an}}
\newcommand{\dd}{\mathrm{d}}
\newcommand{\ii}{\mathrm{i}}
\newcommand{\dbar}{\bar\partial}
\DeclareMathOperator{\CH}{CH}
\DeclareMathOperator{\Hdg}{Hdg}
\DeclareMathOperator{\cl}{cl}
\DeclareMathOperator{\PD}{PD}
\DeclareMathOperator{\im}{im}
\DeclareMathOperator{\Pic}{Pic}
\DeclareMathOperator{\Hom}{Hom}

% ----------------------------------------------------------------------
% Document information
% ----------------------------------------------------------------------
\title[Hodge Classes]
{Notes on Hodge Classes and Algebraic Cycles}
\author{Jhune Yim} % TODO: name, address, email
\email{jhune\_yim@protonmail.com}
\date{12 November 2026}
\subjclass[2020]{Primary 14C30; Secondary 14C25, 14F25, 32J25}
\keywords{Hodge classes, algebraic cycles, cycle-class map, Lefschetz
  theorems, quintic threefold}

\begin{document}
\raggedbottom

\begin{abstract}
This is an expository account of the Hodge conjecture predicts that
every rational cohomology class of type $(p,p)$ on a smooth complex
projective variety is the class of a rational cycle of codimension
$p$. We construct the cycle-class map from fundamental classes and
explain the Hodge-type condition. The analytic exponential sequence
and GAGA give the divisor case, and Hard Lefschetz gives the curve
case. On a smooth quintic threefold the Lefschetz argument represents the
generator of $H^4$ only with denominator five, while a line represents it
integrally. Examples separate the complex Hodge summands from the rational
classes they contain. The final sections discuss integral and K\"ahler
counterexamples, variation in families, and recent results on abelian
varieties.
\end{abstract}

\maketitle

\setcounter{tocdepth}{1}
\tableofcontents

\section{The conjecture}\label{sec:question}

A smooth complex projective variety can be studied through its polynomial
equations, through the topology of its complex points, and through its complex
analytic structure. These descriptions retain different information. The Hodge
conjecture proposes an exact criterion for deciding which rational cohomology
classes arise from algebraic subvarieties.

Fix a connected smooth projective variety $X$ over $\C$ of complex dimension $n$.
Unless a distinction matters, write $X$ also for its associated complex manifold
$X^{\an}$. For $0\leq p\leq n$, define
\begin{equation}\label{eq:hodgeclasses}
 \Hdg^p(X):=\{\alpha\in H^{2p}(X,\Q):
                 \alpha_{\C}\in H^{p,p}(X)\}.
\end{equation}
Here $\alpha_{\C}$ is the image under extension of coefficients from $\Q$ to
$\C$. Equivalently, $\Hdg^p(X)=H^{2p}(X,\Q)\cap H^{p,p}(X)$, with the
intersection taken inside $H^{2p}(X,\C)$.

\begin{conjecture}[Rational Hodge conjecture]\label{conj:hodge}
For every $\alpha\in\Hdg^p(X)$, there exist finitely many irreducible closed
algebraic subvarieties $Z_i\subseteq X$ of complex codimension $p$ and rational
numbers $q_i$ such that
\begin{equation}\label{eq:main}
 \alpha=\sum_i q_i\,\cl(Z_i).
\end{equation}
\end{conjecture}

This is the standard rational formulation \cite{delignecmi2006,cmihodge}. The equality
takes place in cohomology. A representing cycle is a formal sum of
subvarieties with rational coefficients of either sign, and its components
are defined over the base field $\C$.

\section{Prerequisites and conventions}\label{sec:prereq}

The statement and the elementary examples need the following background.
From linear algebra we use vector spaces over $\Q$ and $\C$, tensor
products, duals, exterior powers, and complex conjugation. These enter
through extension of coefficients and the decomposition of cohomology. From
algebraic geometry we use projective varieties, irreducibility, dimension
and codimension, and smoothness. From topology we use singular homology and
cohomology, the cup product, and Poincar\'e duality on a compact oriented
manifold \cite{hatcher2002}. From complex geometry we use differential forms and
their types, Hermitian and K\"ahler metrics, and the Hodge decomposition
\cite{griffithsharris1978,voisin2002hodge,voisin2003hodge2}. The cycle theory, meaning rational equivalence,
Chow groups, divisors and line bundles, and first Chern classes, is
developed by Fulton \cite{fulton1998}. Schemes, derived categories, and motives
are not needed for the statement. The later sections on further directions
introduce additional notions as they arise.

All dimensions and codimensions of algebraic varieties below are complex
unless stated otherwise, so $\dim_{\R}X^{\an}=2n$. For a coefficient ring
$R=\Z,\Q,\R,\C$, the notation $H^k(X,R)$ refers to singular cohomology
of $X^{\an}$. With a sheaf as coefficient, we specify the algebraic or
analytic space whenever both are in use.

\section{Projective varieties and their analytic geometry}

\subsection{Projective space and polynomial equations}

The complex points of projective space have the description
\[
 \PP^N(\C)=(\C^{N+1}\setminus\{0\})/\C^*,
 \qquad
 [z_0:\cdots:z_N]=[\lambda z_0:\cdots:\lambda z_N]
 \quad(\lambda\neq0).
\]
A projective algebraic variety admits a closed embedding into the scheme
$\PP^N_{\C}$. In such an embedding it is given by homogeneous equations
\[
 F_1(z_0,\ldots,z_N)=\cdots=F_r(z_0,\ldots,z_N)=0.
\]
Homogeneity ensures that the equations make sense on projective points.
For example, $z_0z_2-z_1^2=0$ defines a smooth conic in $\PP^2_{\C}$.

Projectivity does not imply smoothness. Smoothness is a separate requirement:
locally the tangent space must have the expected dimension. For a
hypersurface, the Jacobian criterion says that its partial derivatives must
not vanish simultaneously at a projective point on the hypersurface. The
general criterion compares Jacobian rank with codimension \cite{griffithsharris1978}.

\subsection{Compactness, complex structure, and the K\"ahler condition}

The analytic space of a smooth projective $X$ is a compact complex manifold.
Compactness follows because complex projective space is compact and $X$ is
closed. Smooth complex manifolds are canonically oriented as real manifolds.
The Fubini--Study metric on projective space restricts to a K\"ahler metric on
$X$.

A Hermitian metric has an associated real differential form $\omega$ of type
$(1,1)$. It is K\"ahler when $\dd\omega=0$. This condition gives the identities
between differential operators that make the Hodge decomposition possible.
Thus the relevant implications are
\[
 \text{smooth projective over }\C
 \ \Longrightarrow\ \text{compact K\"ahler}
 \ \Longrightarrow\ \text{compact complex manifold}.
\]
The projective embedding provides the integral K\"ahler class
$c_1(\OO_X(1))$, a fact we use when applying Hard Lefschetz over $\Q$
\cite{voisin2002hodge,gaganotes}.

\section{Homology, cohomology, and geometric classes}\label{sec:cohomology}

\subsection{Chains, cycles, and cochains}

Singular homology is built from continuous maps of simplices into a space.
With coefficients in a ring $R$, the singular chain groups $C_k(X,R)$ have
boundary maps $\partial_k$ satisfying $\partial_k\partial_{k+1}=0$. Their
homology is
\[
 H_k(X,R)=\ker(\partial_k)/\im(\partial_{k+1}).
\]
A \emph{topological cycle} is a chain with zero boundary. This usage of
``cycle'' is related to, but different from, an algebraic cycle.

Cochains are $R$-linear functions on chains, with coboundary induced by
$\partial$. Their cohomology is
\[
 H^k(X,R)=\ker(\delta^k)/\im(\delta^{k-1}).
\]
Over $\Q$, the evaluation pairing gives a natural identification
\[
 H^k(X,\Q)\cong\Hom_{\Q}(H_k(X,\Q),\Q).
\]
Over $\Z$, the universal coefficient theorem gives
$\operatorname{Ext}^1_{\Z}(H_{k-1}(X,\Z),\Z)$ as the kernel of evaluation
onto $\Hom_{\Z}(H_k(X,\Z),\Z)$.

The intuition of detecting ``holes'' is helpful for first examples, but
cohomology also records orientation and intersection data. Its multiplication,
the cup product,
\[
 H^a(X,R)\times H^b(X,R)\longrightarrow H^{a+b}(X,R),
 \qquad (u,v)\longmapsto u\smile v,
\]
is graded commutative: $u\smile v=(-1)^{ab}v\smile u$. The ring structure
carries more information than the separate groups \cite{hatcher2002}.

\subsection{de Rham cohomology}

Write $A^k(X,\C)$ for the smooth complex-valued $k$-forms on the manifold.
They form the de Rham complex, with cohomology
\[
 H^k_{\mathrm{dR}}(X,\C)
 =\frac{\ker(\dd:A^k(X,\C)\to A^{k+1}(X,\C))}
 {\im(\dd:A^{k-1}(X,\C)\to A^k(X,\C))}.
\]
The de Rham theorem identifies it with $H^k(X,\C)$, compatibly with
integration and products. A cohomology class is an equivalence class of
closed forms modulo exact forms. It has no distinguished representative
until a metric is chosen.

Rationality is extra information. A complex de Rham class comes from
$H^k(X,\Q)$ exactly when its integrals over integral $k$-cycles are rational
\cite{griffithsharris1978}. Torsion cycles have zero periods, so they impose no condition.

\subsection{Fundamental classes and the degree shift}

For the compact oriented real $2n$-manifold $X$, Poincar\'e duality gives
\[
 \PD:H_{2n-k}(X,\Q)\xrightarrow{\sim}H^k(X,\Q).
\]
If $Z\subseteq X$ is a smooth closed subvariety of complex codimension $p$,
then its real dimension is $2n-2p$. Its fundamental homology class has a
Poincar\'e dual
\begin{equation}\label{eq:cyclePD}
 \cl(Z)=\PD([Z]_{\mathrm{hom}})\in H^{2p}(X,\Q).
\end{equation}
This is the reason that codimension $p$ corresponds to degree $2p$.

The construction extends to singular irreducible subvarieties, which also
have fundamental homology classes. Smoothness of $X$ does not require
smoothness of the $Z_i$. One characterization is
\begin{equation}\label{eq:integration}
 \int_X \cl(Z)\smile[\eta]=\int_Z\eta
\end{equation}
for closed forms $\eta$ of real degree $2n-2p$, interpreted through the
fundamental class, or integration over the regular locus, if $Z$ is singular.
We use $[Z]$ for the cycle and $\cl(Z)$ for its cohomology class \cite{fulton1998,voisin2003hodge2}.

\section{Hodge decomposition and rational Hodge classes}\label{sec:hodge}

\subsection{Forms of type \texorpdfstring{$(a,b)$}{(a,b)}}

In holomorphic coordinates $z_1,\ldots,z_n$, a form of type $(a,b)$ is a sum
of terms
\[
 f_{I,J}(z,\bar z)\,
 \dd z_{i_1}\wedge\cdots\wedge\dd z_{i_a}
 \wedge\dd\bar z_{j_1}\wedge\cdots\wedge\dd\bar z_{j_b}.
\]
Its total degree is $a+b$. The exterior derivative decomposes as
$\dd=\partial+\dbar$, with bidegrees $(1,0)$ and $(0,1)$.

On a compact Riemannian manifold, Hodge theory assigns a unique harmonic
representative to every de Rham cohomology class, once the metric is fixed.
Harmonicity means annihilation by
$\Delta_{\dd}=\dd\dd^*+\dd^*\dd$. On a compact K\"ahler manifold the
K\"ahler identities give $\Delta_{\dd}=2\Delta_{\dbar}$, so harmonic forms
split by type \cite{voisin2002hodge}. The companion notes derive the identities in
holomorphic normal coordinates and give the harmonic proof of the
decomposition \cite[\S\S5--6]{gaganotes}.

\begin{theorem}[Hodge decomposition]
For a compact K\"ahler manifold,
\begin{equation}\label{eq:decomposition}
 H^k(X,\C)=\bigoplus_{a+b=k}H^{a,b}(X),
 \qquad \overline{H^{a,b}(X)}=H^{b,a}(X).
\end{equation}
The space $H^{a,b}(X)$ consists of de Rham classes admitting a closed
representative of pure type $(a,b)$ and is naturally identified with
Dolbeault cohomology $H^b(X,\Omega_X^a)$.
\end{theorem}

Although harmonic representatives depend on the K\"ahler metric, this
decomposition depends on the complex structure and is independent of the
chosen K\"ahler metric. The dimensions $h^{a,b}=\dim_{\C}H^{a,b}(X)$ are
Hodge numbers, with $b_k=\sum_{a+b=k}h^{a,b}$.

\subsection{The rational structure}

Extension of coefficients gives
\[
 H^k(X,\Q)\otimes_{\Q}\C\cong H^k(X,\C).
\]
However, the individual Hodge summands need not be defined over $\Q$.
In particular, $H^{p,p}(X)$ is a complex vector space, whereas
$\Hdg^p(X)$ is the rational subspace specified by
\eqref{eq:hodgeclasses}. Knowing $h^{p,p}$ alone does not determine the
dimension of $\Hdg^p(X)$.

The linear algebra is already visible for $\Q^2\subset\C^2$: a complex line
of irrational slope contains no nonzero rational point.

The Hodge filtration is
\[
 F^rH^k(X,\C)=\bigoplus_{a\geq r}H^{a,k-a}(X).
\]
Because a rational class is fixed by complex conjugation,
\begin{equation}\label{eq:filtration}
 \Hdg^p(X)=H^{2p}(X,\Q)\cap F^pH^{2p}(X,\C).
\end{equation}
Indeed, membership in $F^p$ and its conjugate forces both Hodge indices to
be at least $p$. Since their sum is $2p$, both equal $p$ \cite{delignecmi2006}.

\subsection{Hodge structures and Tate twists}

A pure rational Hodge structure of weight $k$ is a finite-dimensional
$\Q$-vector space $V$ with a decomposition
$V_{\C}=\bigoplus_{a+b=k}V^{a,b}$ satisfying
$\overline{V^{a,b}}=V^{b,a}$. Cohomology $H^k(X,\Q)$ is the principal
example here. A morphism of Hodge structures is a rational linear map
preserving these summands.

The Tate Hodge structure $\Q(1)$ has weight $-2$ and type $(-1,-1)$.
The twist
\[
 H^{2p}(X,\Q(p)):=H^{2p}(X,\Q)\otimes\Q(p)
\]
has weight zero, so a Hodge class becomes a class of type $(0,0)$.
Equivalently,
\[
 \Hdg^p(X)\cong
 \Hom_{\mathrm{HS}}\bigl(\Q(0),H^{2p}(X,\Q(p))\bigr).
\]
When compatible integral or de Rham normalizations are included, Tate twists
carry conventional factors of $2\pi\ii$. Equation~\eqref{eq:main} uses
the untwisted Betti cycle class \cite{deligne1982hodge}.

\section{Sheaves, line bundles, and Chern classes}\label{sec:bundles}

A \emph{sheaf} assigns data to open sets with restriction maps, so
compatible local data glue uniquely. On $X^{\an}$, the sheaf $\OO_{X^{\an}}$
consists of holomorphic functions, its units $\OO_{X^{\an}}^*$ are the
nowhere-vanishing holomorphic functions, and $\Omega^a_{X^{\an}}$
consists of holomorphic $a$-forms. Their algebraic counterparts are
defined on Zariski open sets. Regular functions there are locally
quotients of polynomials with nonvanishing denominator. Rational
functions are specified on a dense open set and may have poles elsewhere.

Sheaf cohomology $H^j(X,\mathcal F)$ measures obstructions to gluing local
data. A short exact sequence of sheaves produces a long exact sequence
of these groups. This is the formal fact used in the divisor proof.
GAGA compares coherent sheaf cohomology on the projective scheme
with that on its analytic space. In particular,
$H^b(X,\Omega^a_{X/\C})\cong H^b(X^{\an},\Omega^a_{X^{\an}})$
\cite{serre1956gaga,gaganotes}. Dolbeault theory computes the analytic group
using smooth forms \cite{voisin2002hodge}:
\[
 H^b(X^{\an},\Omega^a_{X^{\an}})\cong
 \frac{\ker(\dbar:A^{a,b}(X)\to A^{a,b+1}(X))}
      {\im(\dbar:A^{a,b-1}(X)\to A^{a,b}(X))}.
\]
The denominator is zero when $b=0$.

A \emph{line bundle} is a locally trivial family of one-dimensional complex
vector spaces over $X$, with holomorphic or algebraic transition
functions as appropriate. Isomorphism classes of line bundles form the
Picard group under tensor product. Holomorphic transition functions identify
$\Pic(X^{\an})$ with $H^1(X^{\an},\OO_{X^{\an}}^*)$. GAGA identifies this
group with the algebraic Picard group $\Pic(X)$.
A \emph{divisor} on smooth $X$ is an integral codimension-one cycle.
It determines a line bundle $\OO_X(D)$. Conversely, a nonzero rational
section of a line bundle gives a divisor of zeros and poles.

For a rank-$r$ complex vector bundle $E$, the Chern classes
$c_j(E)\in H^{2j}(X,\Z)$ are characteristic classes, with
$c_0(E)=1$ and $c_j(E)=0$ for $j>r$. They are natural under pullback
and satisfy the Whitney product formula for short exact sequences.
For line bundles,
\[
 c_1(L\otimes M)=c_1(L)+c_1(M),\qquad
 c_1(\OO_X(D))=\cl(D).
\]
An algebraic vector bundle on smooth $X$ also has Chern classes in
$\CH^j(X)$, whose images under the cycle-class map are these cohomological
Chern classes \cite{fulton1998}.

On the spaces considered here, a coherent sheaf of modules over the
structure sheaf locally has a finite presentation. Sheaves of sections of
vector bundles are coherent. A coherent sheaf can also be supported on a
proper subvariety. On a smooth projective variety, coherent sheaves have finite
locally free resolutions, so one can define their Chern classes through
such resolutions. A global resolution of this kind need not exist on a
general compact K\"ahler manifold, a distinction relevant to the
counterexamples discussed later \cite{voisin2002kahler}.

An ample line bundle is one whose sufficiently high positive tensor
power gives a projective embedding. Its first Chern class is an
ample divisor class. The hyperplane bundle $\OO_{\PP^N}(1)$ is the
basic example. Its sections are linear homogeneous equations, and the
class of a hyperplane is its first Chern class.

\section{Algebraic cycles and the cycle-class map}\label{sec:cycles}

\subsection{Formal cycles, coefficients, and equivalence}

\begin{definition}
The group of codimension-$p$ algebraic cycles is the free abelian group
\[
 Z^p(X)=\bigoplus_{\substack{Z\subseteq X\ \mathrm{irreducible,\ closed}\\
                              \operatorname{codim}_X Z=p}}\Z\,[Z].
\]
An element is a finite sum $z=\sum_i m_i[Z_i]$ with $m_i\in\Z$.
A cycle with rational coefficients is an element of
$Z^p(X)_{\Q}:=Z^p(X)\otimes_{\Z}\Q$.
\end{definition}

A component $Z_i$ of a codimension-$p$ cycle has dimension $n-p$.
Divisors are codimension-one cycles. Zero-cycles are finite combinations
of points and have codimension $n$. A cycle is \emph{effective} if its
coefficients are nonnegative. The Hodge conjecture imposes no effectiveness
requirement.

Rational equivalence is generated by divisors of nonzero rational functions
on irreducible subvarieties $W\subseteq X$ of codimension $p-1$, pushed
forward to $X$. If $W$ is singular, these divisors may be defined using its
normalization and pushforward. The quotient is the Chow group
\[
 \CH^p(X)=Z^p(X)/\!\sim_{\mathrm{rat}},
 \qquad \CH^p(X)_{\Q}=\CH^p(X)\otimes\Q.
\]
The adjective ``rational'' refers to the rational functions used in the
definition. Rationally equivalent cycles have the same cohomology class, and
the kernel of the cycle-class map on $\CH^p(X)$ can still be nonzero. Fulton
gives the formal construction and its functorial properties \cite{fulton1998}.

\subsection{Why every algebraic cycle class is a Hodge class}

\begin{proposition}\label{prop:necessary}
An irreducible closed algebraic subvariety $Z\subseteq X$ of complex
codimension $p$ determines an integral class whose complexification has type
$(p,p)$. In particular, its rational class belongs to $\Hdg^p(X)$.
\end{proposition}

\begin{proof}[Geometric explanation]
First suppose $Z$ is smooth. It has complex dimension $n-p$, so integration
over $Z$ can detect a form of total degree $2n-2p$ only through its
$(n-p,n-p)$ component. The cup product respects bidegree, so Poincar\'e
duality pairs each Hodge summand perfectly with the complementary one.
It follows that $\cl(Z)$ has only a $(p,p)$ component. For singular $Z$,
integration over its regular locus gives a closed current of bidegree
$(p,p)$ representing the same fundamental class. The fundamental class
is integral, which also gives the asserted rationality.
\end{proof}

A current is a continuous linear functional on smooth test forms.
In this sense integration over a subvariety defines a generalized
differential form \cite{griffithsharris1978,voisin2003hodge2}.

By linearity and invariance under rational equivalence, we obtain
\begin{equation}\label{eq:cyclemap}
 \cl^p_{\Q}:\CH^p(X)_{\Q}\longrightarrow\Hdg^p(X),
 \qquad [z]\longmapsto\cl(z).
\end{equation}
The Hodge conjecture is exactly the assertion that this map is surjective
for every $p$. Its image is also the image of the map from
$Z^p(X)_{\Q}$. Passing to Chow groups lets us study its kernel
modulo rational equivalence (Section~\ref{subsec:abel-jacobi}).

\section{Equivalent formulations}

Let
\[
 H^{2p}_{\mathrm{alg}}(X,\Q):=
 \operatorname{span}_{\Q}\{\cl(Z):\operatorname{codim}_X Z=p\}.
\]
Proposition~\ref{prop:necessary} proves the inclusion
\[
 H^{2p}_{\mathrm{alg}}(X,\Q)\subseteq\Hdg^p(X).
\]
The conjecture predicts equality.

There is also a useful denominator-clearing formulation. If
$\alpha\in\Hdg^p(X)$, equation~\eqref{eq:main} is equivalent to the
existence of an integer $N>0$ and an integral algebraic cycle $z$ such that
\begin{equation}\label{eq:multiple}
 N\alpha=\cl(z)\quad\text{in }H^{2p}(X,\Q).
\end{equation}
For an integral Hodge class $\beta$, the conjecture likewise implies that
some positive multiple of $\beta$ is an integral algebraic cycle class,
with equality even in integral cohomology after multiplying again to kill
any torsion discrepancy. It does not imply that $\beta$ itself is algebraic
integrally \cite{delignecmi2006}.

Some points of interpretation follow directly from the statement. The
conjecture produces an algebraic representative of a class. It says nothing
about the other representatives of the same class, by differential forms or
by topological cycles. The coefficients may be negative, so a Hodge class
need not be the class of a single subvariety or of an effective cycle. The
statement is one of existence: it gives no procedure for finding equations
of the subvarieties and no bound on their degrees. Equality of cycle classes
is weaker than rational equivalence, so the conjecture describes the image
of the cycle-class map and leaves its kernel untouched.

\section{Examples}\label{sec:examples}

\subsection{Projective space}

Let $h=c_1(\OO_{\PP^n}(1))$ be the hyperplane class. Then
\[
 H^*(\PP^n,\Q)=\Q[h]/(h^{n+1}),\qquad \deg h=2.
\]
A linear subspace $\PP^{n-p}\subset\PP^n$ has class $h^p$, so
\[
 \Hdg^p(\PP^n)=H^{2p}(\PP^n,\Q)=\Q h^p.
\]
Every $qh^p$ is represented by the rational cycle $q[\PP^{n-p}]$.
So the conjecture holds for projective space, because the cohomology ring
has explicit algebraic generators \cite{hatcher2002,griffithsharris1978}.

\subsection{Plane curves}

In $\PP^2_{\C}$, a projective curve $C$ of degree $d$ is a divisor with
$\OO(C)\cong\OO(d)$. Consequently
\[
 \cl(C)=c_1(\OO(d))=d h.
\]
The smooth conic $z_0z_2-z_1^2=0$ has class $2h$. The union of two
distinct lines, for example $z_0z_1=0$, also has class $2h$. The first
curve is irreducible and smooth. The second is reducible and singular at
the intersection point. Their common class records their degree.

For a line $L$, the class $\tfrac12h$ is represented by the rational cycle
$\tfrac12[L]$. It cannot be the class of an integral divisor on $\PP^2$,
since those have integral degree. The class $-h$ is a Hodge class but
cannot be represented by an effective divisor, whose degree is nonnegative.

\subsection{An elliptic curve}

An elliptic curve over $\C$ is a smooth projective curve of genus one
with a chosen point. For such a curve $E$,
\[
 H^1(E,\C)=H^{1,0}(E)\oplus H^{0,1}(E),
 \qquad h^{1,0}=h^{0,1}=1.
\]
This nonzero odd-degree cohomology cannot be the class of an algebraic
subvariety: cycle classes have even degree. Meanwhile $H^2(E,\Q)$ is
generated by the class of a point. The conjecture applies to the latter
and makes no assertion that the former is generated by algebraic cycle
classes.

In even degree the same exclusion applies to type. By
Proposition~\ref{prop:necessary}, a rational class of degree $2p$ whose
complexification has a component outside $H^{p,p}$ is not algebraic.

\subsection{K3 surfaces}

A complex K3 surface is a compact simply connected complex surface with
trivial canonical bundle. For a K3 surface $S$ one has $h^{1,1}(S)=20$. By
Theorem~\ref{thm:11} below, the rational Hodge classes in degree two are
the rational span of divisor classes. Its dimension is the Picard number
$\rho(S)$, which takes every value from $1$ to $20$ among projective K3
surfaces. So $h^{1,1}=20$ does not give twenty rationally independent
divisor classes. What varies is the position of the rational lattice inside
complex cohomology \cite{voisin2002hodge,huybrechts2016k3}.

\section{Known cases}\label{sec:known}

\subsection{The Lefschetz \texorpdfstring{$(1,1)$}{(1,1)} theorem}

\begin{theorem}[Lefschetz $(1,1)$ theorem]\label{thm:11}
On a smooth complex projective variety, every integral cohomology class
whose complexification has type $(1,1)$ is the first Chern class of an
algebraic line bundle, and hence the class of an integral divisor.
Consequently the rational Hodge conjecture holds for $p=1$.
\end{theorem}

\begin{proof}
Put $M=X^{\an}$. The analytic exponential sequence
\[
 0\longrightarrow\Z\longrightarrow\OO_M
 \xrightarrow{\ f\mapsto\exp(2\pi\ii f)\ }\OO_M^*
 \longrightarrow1
\]
gives the exact sequence
\[
 \Pic(M)\xrightarrow{\ c_1\ }H^2(M,\Z)\longrightarrow H^2(M,\OO_M).
\]
Dolbeault theory identifies the last group with $H^{0,2}(M)$, and the map
from integral cohomology is complexification followed by projection to this
summand. Indeed, projection of a smooth form onto its $(0,\bullet)$
component induces the map from the de Rham complex to the Dolbeault
resolution of $\OO_M$. An integral class of type $(1,1)$ has zero $(0,2)$
component, so by exactness it is the first Chern class of a holomorphic
line bundle. Conversely, the complexification of an integral class is real,
so vanishing of its $(0,2)$ component forces vanishing of its $(2,0)$
component. The image of $c_1$ therefore consists exactly of the integral
classes of type $(1,1)$.

GAGA algebraizes the holomorphic line bundle on projective $X$
\cite{serre1956gaga,gaganotes}. A nonzero rational section of the resulting
algebraic line bundle $L$ gives a divisor $D$ with
$c_1(L)=\cl(D)$. Multiplying a rational $(1,1)$ class by a suitable
positive integer reduces the rational assertion to this integral case.
\end{proof}

The theorem covers integral torsion classes as well: their complexification
is zero, and exactness places them in the image of $c_1$. Nothing comparable
holds in higher codimension. No analogue of the exponential sequence
classifies cycles there, and the integral statement fails
(Section~\ref{sec:limits}) \cite{voisin2002hodge}.

\subsection{Curve classes and Hard Lefschetz}

Let $h$ be the class of an ample divisor and write $L(u)=h\smile u$.
A positive multiple of $h$ is a hyperplane class for a projective
embedding, so $h$ is represented by a K\"ahler form. Hard Lefschetz
over $\C$ is proved in the companion notes by the action on harmonic
forms \cite[\S7]{gaganotes}. Cup product by $h$ is defined over $\Q$,
so the map that becomes an isomorphism after extension to $\C$ is
already an isomorphism over $\Q$:
\begin{equation}\label{eq:HL}
 L^{n-k}:H^k(X,\Q)\xrightarrow{\sim}H^{2n-k}(X,\Q)
 \qquad(0\leq k\leq n).
\end{equation}
Since $h$ has type $(1,1)$, for $n\geq2$ the isomorphism
\[
 L^{n-2}:H^2(X,\Q)\xrightarrow{\sim}H^{2n-2}(X,\Q)
\]
takes type $(1,1)$ to type $(n-1,n-1)$, and its inverse takes rational
Hodge classes to rational $(1,1)$ classes. By Theorem~\ref{thm:11},
the preimage of such a class is a rational divisor class. Its product with
$h^{n-2}$ is algebraic, because the cycle-class map respects intersection
products. Hence the rational Hodge conjecture holds for codimension
$n-1$, that is, for one-dimensional cycles \cite{voisin2003hodge2}.

\begin{example}[Curve classes on a smooth quintic]
\label{ex:quintic}
Let $X\subset\PP^4_{\C}$ be a smooth quintic threefold and
$h=c_1(\OO_X(1))$. The companion notes compute
$H^4(X,\Q)=\Q h^2$, with $h^2$ of type $(2,2)$, and
$\int_X h^3=5$ \cite[\S8]{gaganotes}.
The curve $C=X\cap H_1\cap H_2$ cut out by two general hyperplanes
has class $\cl(C)=h^2$. Thus every rational Hodge class of degree four
on $X$ is a rational multiple of $\cl(C)$.

The integral picture is finer. The homological form of the weak Lefschetz
theorem and Poincar\'e duality give
$H^4(X,\Z)\cong H_2(X,\Z)\cong\Z$ \cite[Ch.~1]{voisin2003hodge2}. Let $\ell$ be the
generator with $\int_X h\smile\ell=1$. Then $h^2=5\ell$, so the map
\[
 h\smile{-}:H^2(X,\Z)\longrightarrow H^4(X,\Z)
\]
has image of index five, and the Lefschetz argument represents $\ell$ only
by the rational cycle $\tfrac15[C]$.

On the quintic the denominator can be removed by other means. The lines on
$X$ are the zeros of a section of the rank-six bundle
$\operatorname{Sym}^5S^\vee$ on the six-dimensional Grassmannian of lines
in $\PP^4$, where $S$ is the tautological subbundle. Its top Chern class is
$2875\neq0$, so $X$ contains a line $R$ \cite[Ch.~6]{eisenbudharris2016}.
Since $\int_X h\smile\cl(R)=\deg R=1$, we get $\cl(R)=\ell$. For
hypersurfaces of other degrees the generator of $H^4(X,\Z)$ need not be
algebraic (Section~\ref{sec:limits}).
\end{example}

\subsection{Consequences in small dimension}

Codimension zero is generated by $[X]$, and top-degree cohomology of a
connected $X$ is generated by a point. Together with the divisor and curve
cases, this proves the rational Hodge conjecture for every smooth complex
projective variety of dimension at most three. The first dimension where
the general problem is not covered by these arguments is $n=4$, in
codimension $p=2$ \cite{voisin2016hodge}.

Some special fourfolds are understood. Zucker proved the rational Hodge
conjecture for smooth cubic fourfolds $X\subset\PP^5_{\C}$, where the
nontrivial case is $H^4(X,\Q)$ \cite{zucker1977}.

\subsection{Abelian varieties}

An abelian variety is a connected projective algebraic group. Over $\C$,
its analytification is a complex torus with a polarization from an ample
line bundle. Its cohomology is the exterior algebra on $H^1$, and its Hodge
classes need not be generated by divisor classes. The basic examples are
the Weil classes, which arise when an imaginary quadratic field acts on an
abelian variety of even dimension with balanced action on the holomorphic
tangent space.

Markman proved that the Weil classes on abelian sixfolds of Weil type of a
fixed discriminant are algebraic, for every imaginary quadratic field,
using secant sheaves and deformation methods. By a degeneration argument of
Schoen this gives the Weil classes on all abelian fourfolds of Weil type,
and the Hodge conjecture for abelian fourfolds follows \cite{markman2025}.
Voisin's Bourbaki expos\'e explains how the same input, with reductions due
to Moonen and Zarhin, gives the Hodge conjecture for complex abelian
varieties of dimension at most five
\cite[Th\'eor\`eme~1.3]{voisin2026bourbaki}.

The general rational Hodge conjecture remains open \cite{cmihodge}.

\section{Significance}\label{sec:importance}

\subsection{A criterion for algebraicity}

The conjecture would identify the algebraic part of rational cohomology
using two pieces of data: the rational structure and the Hodge
decomposition. The implication already known is
\[
 \text{algebraic cycle class}\ \Longrightarrow\
 \text{rational class of type }(p,p).
\]
Its conjectural converse is an existence statement. A class found by
analytic or cohomological reasoning would force algebraic subvarieties to
exist before any equations for them were known \cite{cmihodge}.

The criterion works at the level of rational cycle classes. It would not
recover the equations of $X$ or single out a cycle, since many cycles share
a class, as the conics of Section~\ref{sec:examples} show.

\subsection{Intersection theory}

For properly intersecting subvarieties, cup products of cycle classes
encode their intersection cycle classes. In complementary dimensions,
integrating the cup product gives an intersection number. For example,
two plane curves of degrees $d$ and $e$ with no common component have
intersection number
\[
 \int_{\PP^2}(dh)\smile(eh)=de,
 \qquad \int_{\PP^2}h^2=1.
\]
The intersection multiplicities are essential when intersections are
not transverse. Refined intersection theory extends these constructions
beyond the naive set-theoretic intersection \cite{fulton1998}.

If the conjecture holds, every rational Hodge class takes part in this
intersection theory as an algebraic cycle class. Cycles with the same
cohomology class, or the same intersection numbers, still need not be
rationally equivalent.

\subsection{Correspondences}

For smooth projective $X,Y$ with $\dim X=n$, a codimension-$r$ cycle
$\Gamma$ on $X\times Y$ acts by
\begin{equation}\label{eq:correspondence}
 \Gamma_*(u)=(\operatorname{pr}_Y)_*
 \bigl(\operatorname{pr}_X^*u\smile\cl(\Gamma)\bigr)
 \in H^{k+2r-2n}(Y,\Q),\quad u\in H^k(X,\Q).
\end{equation}
Pullback, cup product, and integration along fibers explain the degree
shift. Such a cycle is an algebraic correspondence. When $r=n$, its
action preserves degree and Hodge type.

By K\"unneth decomposition and Poincar\'e duality, a rational morphism
of Hodge structures $H^k(X,\Q)\to H^k(Y,\Q)$ determines a rational
$(n,n)$-class in the appropriate summand of $H^{2n}(X\times Y,\Q)$.
The Hodge conjecture for the product would make that class algebraic,
realizing the map by a correspondence. This is why the conjecture for
products matters: it would give a geometric origin to every morphism of
Hodge structures between the cohomology groups of two varieties
\cite{voisin2025coniveau}.

Buskin's theorem is an instance. Every rational Hodge isometry between the
degree-two cohomologies of two projective K3 surfaces is induced by an
algebraic correspondence \cite{buskin2019}. An isometry here is a map
preserving the intersection form.

\subsection{Motives and the Tate conjecture}

The theory of motives seeks a framework in which varieties are compared
using algebraic correspondences and different cohomology theories appear
as realizations of common geometric objects. The Hodge conjecture would
ensure that certain Hodge-theoretic relations come from actual
correspondences, one of the properties sought in a theory of motives \cite{delignecmi2006}.

Over arithmetic fields, the Tate conjecture asks for an analogous
algebraicity criterion using Galois-invariant $\ell$-adic cohomology
classes, with a Tate twist. For a variety over a finite field, suitable
Frobenius-invariant classes play the role occupied here by Hodge classes.
Both conjectures ask which cohomology classes have an algebraic origin.
They are different statements, and neither is known to imply the other in
general \cite{tate1965,milne2008tate}.

\section{Integral, K\"ahler, and singular variants}\label{sec:limits}

\subsection{Integral coefficients}

For integral coefficients, define
\[
 \Hdg^p(X,\Z):=
 \{\beta\in H^{2p}(X,\Z):\beta_{\C}\in H^{p,p}(X)\}.
\]
This preimage definition is necessary because the map from integral to
complex cohomology kills torsion and need not be injective. Every torsion
class belongs to this group, since its complexification is zero.

One might ask whether $\CH^p(X)\to\Hdg^p(X,\Z)$ is always surjective.
Atiyah and Hirzebruch showed that the answer is negative: algebraic cycle
classes satisfy additional topological constraints, invisible to the
Hodge-type condition, and some integral Hodge classes violate them
\cite{atiyahhirzebruch1962}. Torsion is a fundamental source of this failure.

Failure is not confined to torsion classes. Let $X\subset\PP^4_{\C}$ be a
very general hypersurface of degree $d$, and let $\ell$ generate
$H^4(X,\Z)\cong\Z$ with $\int_X h\smile\ell=1$, as in
Example~\ref{ex:quintic}, so that $h^2=d\,\ell$. Koll\'ar showed that if
$p^3$ divides $d$ for a prime $p\geq5$, then every curve on $X$ has degree
divisible by $p$. In that case $\ell$ is a non-torsion integral Hodge class
that is not algebraic, although $d\,\ell=\cl(C)$ is, for a plane section
$C$ \cite{voisin2006integral}. The rational conjecture is unaffected, since
$\ell=\tfrac1d\cl(C)$ in $H^4(X,\Q)$.

\subsection{Compact K\"ahler manifolds}

Compact K\"ahler manifolds still have Hodge decomposition, so it is
natural to ask for an analogous algebraicity statement there. In that
setting one must replace algebraic subvarieties by closed complex
analytic subvarieties, or consider Chern classes of suitable analytic
objects. Unrestricted extensions of this kind fail. Voisin constructs
compact K\"ahler examples where even rational Hodge classes are not
generated by the Chern classes of coherent analytic sheaves
\cite{voisin2002kahler}.

On a projective variety every closed analytic subvariety is algebraic by
Chow's theorem \cite{chow1949,gaganotes}, so there the analytic and algebraic
forms of the question coincide.

\subsection{Singular and open varieties}

For singular or noncompact algebraic varieties, ordinary cohomology has
in general a mixed Hodge structure in place of the pure structure above.
A mixed Hodge structure has an increasing rational weight filtration
$W_\bullet$ and a decreasing complex Hodge filtration $F^\bullet$,
such that each graded piece $\operatorname{Gr}^W_k$ is a pure Hodge
structure of weight $k$.
Poincar\'e duality also requires modification. Analogues of the conjecture
can be formulated with mixed Hodge theory, homology, or intersection
cohomology, each with its own hypotheses \cite{voisin2003hodge2,voisin2025coniveau}.

\section{Further directions}\label{sec:research}

\subsection{Why Hodge theory does not produce cycles}

Hodge theory decomposes a finite-dimensional vector space and describes how
the decomposition varies. Finding a cycle means finding closed
subvarieties, with whatever singularities, degrees, and deformation
behaviour they have. Nothing turns a harmonic $(p,p)$-form into the zero
set of polynomial equations.

The codimension-one proof works because line bundles are controlled by
the exponential sequence. In higher codimension, coherent sheaves and
their Chern classes, deformation theory, and normal functions are partial
substitutes, and none of them produces every Hodge class
\cite{voisin2016hodge}.

\subsection{Hodge loci and variation in families}

Let $f:\mathcal X\to S$ be a smooth projective family over a smooth
complex algebraic base. Locally, the groups $H^{2p}(X_s,\Q)$ form a
constant system of vector spaces, while the Hodge filtration varies
holomorphically. Parallel transport can move a rational class between
nearby fibers, but the class need not stay of type $(p,p)$.

The Hodge locus records where such classes acquire or retain the required
type. Cattani, Deligne, and Kaplan proved that the Hodge locus in the
base is a countable union of closed algebraic subvarieties
\cite{cdk1995}. Monodromy must be accounted for when tracking one class
globally. It controls whether a locally fixed class defines a global
section of the cohomology local system.

This is what the Hodge conjecture predicts. The locus where a class is
algebraic is a countable union of algebraic subvarieties, by the theory of
Hilbert schemes. The theorem is unconditional evidence for the conjecture,
though it constructs no cycles.

\subsection{Homologically trivial cycles and Abel--Jacobi theory}
\label{subsec:abel-jacobi}

Write
\[
 \CH^p(X)_{\Q,\mathrm{hom}}
 :=\ker\bigl(\cl^p_{\Q}:\CH^p(X)_{\Q}\to H^{2p}(X,\Q)\bigr).
\]
Elements of this group have zero rational cohomology class but need not
vanish modulo rational equivalence. Abel--Jacobi maps and intermediate
Jacobians detect some information discarded by the cycle-class map.
At the integral level, the intermediate Jacobian relevant to codimension
$p$ is the complex torus
\[
 J^p(X)=\frac{H^{2p-1}(X,\C)}
 {F^pH^{2p-1}(X,\C)+\im H^{2p-1}(X,\Z)}.
\]
The integral Abel--Jacobi map is defined on cycles with vanishing integral
cohomology class, modulo rational equivalence. Tensoring the domain
and target with $\Q$ gives the corresponding rational map.
The Abel--Jacobi map refines the information about the kernel. It is
independent of the conjectured surjectivity in degree $2p$
\cite{voisin2003hodge2}.

\subsection{Absolute Hodge classes}

Absolute Hodge theory asks whether a Hodge class remains a Hodge class
under all field automorphisms of $\C$, using the appropriate de Rham
comparison to transport it to the conjugate variety. Algebraic cycles
give absolute Hodge classes. Deligne proved that Hodge classes on
complex abelian varieties are absolute Hodge \cite{deligne1982hodge}.

Absolute Hodge is weaker than algebraic. Deligne's theorem shows that Hodge
classes on abelian varieties behave under automorphisms of $\C$ as algebraic
classes would. It leaves the existence of cycles open.

\subsection{The generalized Hodge conjecture}

The ordinary conjecture concerns degree-$2p$ classes supported, via cycle
classes, on algebraic sets of codimension $p$. A broader question asks
which parts of $H^k(X,\Q)$ can be supported on subsets of prescribed
codimension. The geometric coniveau is the following sum over closed
algebraic subsets $Z\subseteq X$:
\[
 N^cH^k(X,\Q)=
 \sum_{\operatorname{codim}_X Z\geq c}
 \ker\bigl(H^k(X,\Q)\to H^k(X\setminus Z,\Q)\bigr).
\]
A rational Hodge substructure has Hodge coniveau at least $c$ when its
complexification has no components of type $(a,b)$ with $a<c$ or $b<c$.
Let $M^cH^k$ be the largest rational Hodge substructure with that property.
The generalized Hodge conjecture predicts
\[
 N^cH^k(X,\Q)=M^cH^k(X,\Q).
\]
The inclusion $N^cH^k\subseteq M^cH^k$ is known, and the converse is the
content of the conjecture. The condition has to be imposed on a rational
Hodge substructure. Hodge's original version, stated for individual
classes, is false, as Grothendieck observed
\cite{grothendieck1969,voisin2025coniveau}.

\appendix
\section{Notation and degree checks}

\begin{center}\small
\begin{tabularx}{\textwidth}{@{}lX@{}}\toprule
Notation & Meaning \\ \midrule
$n$ & Complex dimension of the smooth projective ambient variety $X$.\\
$p$ & Complex codimension of a cycle with components of dimension $n-p$.\\
$X^{\an}$ & The associated compact complex manifold, of real dimension $2n$.\\
$H^k(X,R)$ & Singular cohomology with coefficient ring or field $R$.\\
$H^{a,b}(X)$ & Hodge summand in complex cohomology of total degree $a+b$.\\
$\Hdg^p(X)$ & Rational classes in $H^{2p}$ whose complexifications have type $(p,p)$.\\
$Z^p(X)_{\Q}$ & Finite formal rational combinations of codimension-$p$ subvarieties.\\
$\CH^p(X)_{\Q}$ & The same cycles modulo rational equivalence.\\
$\cl(Z)$ & The cohomology class of a subvariety, in degree $2p$.\\
$h$ & An ample divisor class, of degree $2$ and Hodge type $(1,1)$.\\
\bottomrule
\end{tabularx}
\end{center}

The two central degree calculations are
\[
 \dim_{\R}Z=2(n-p),\qquad
 \deg\PD([Z]_{\mathrm{hom}})=2n-2(n-p)=2p,
\]
and, for curve classes,
\[
 2+2(n-2)=2n-2,\qquad
 (1,1)+(n-2,n-2)=(n-1,n-1).
\]



\bibliographystyle{amsalpha}
\bibliography{references}
\end{document}

%%% Local Variables:
%%% mode: LaTeX
%%% TeX-master: t
%%% End:
